\section{Ethylic alcohol và độ cồn}

\subsection{Cấu tạo và tính chất vật lí}

Ethylic alcohol có công thức phân tử \ce{C2H6O} và công thức cấu tạo thu gọn
\[
\ce{CH3-CH2-OH}.
\]
Nhóm \ce{-OH} là đặc điểm cấu tạo quan trọng của ethylic alcohol.

Ở điều kiện thường, ethylic alcohol là chất lỏng không màu, có mùi đặc trưng, tan vô hạn trong nước, nhẹ hơn nước và dễ bay hơi. Ethylic alcohol còn là chất dễ cháy.

\subsection{Độ cồn}

Độ cồn cho biết số mililít ethylic alcohol nguyên chất có trong \SI{100}{\milli\liter} dung dịch rượu ở điều kiện quy ước.

\begin{ghinho}{Ý nghĩa của độ cồn}
Rượu $x^\circ$ có $x$ mL ethylic alcohol nguyên chất trong \SI{100}{\milli\liter} dung dịch.
\[
V_{\ce{C2H5OH}}=V_{\text{dd}}\frac{x}{100}.
\]
Nếu cần đổi từ thể tích ethylic alcohol sang khối lượng, dùng $m=DV$.
\end{ghinho}

\begin{vidu}
Một chai chứa \SI{500}{\milli\liter} rượu $40^\circ$. Tính thể tích ethylic alcohol nguyên chất có trong chai.
\end{vidu}
\begin{loigiai}
\[
V_{\ce{C2H5OH}}=500\times\frac{40}{100}=\SI{200}{\milli\liter}.
\]

\textbf{Vậy,} chai rượu chứa \SI{200}{\milli\liter} ethylic alcohol nguyên chất.
\end{loigiai}

\subsection{Tính chất hóa học}

\textbf{Phản ứng cháy}
\[
\ce{C2H5OH + 3O2 -> 2CO2 + 3H2O}.
\]

\textbf{Phản ứng với sodium}
\[
\ce{2C2H5OH + 2Na -> 2C2H5ONa + H2 ^}.
\]
Khi cho sodium vào ethylic alcohol, sodium tan dần và có khí thoát ra.

\begin{vidu}
Cho \SI{4.60}{\gram} ethylic alcohol nguyên chất tác dụng hoàn toàn với sodium dư. Tính thể tích \ce{H2} thu được ở điều kiện có $V_m=\SI{24.8}{\liter\per\mole}$.
\end{vidu}
\begin{loigiai}
\[
n_{\ce{C2H5OH}}=\frac{4.60}{46}=\SI{0.100}{\mole}.
\]

\bangmol{7}
{\ce{2C2H5OH} & + & \ce{2Na} & \rightarrow & \ce{2C2H5ONa} & + & \ce{H2}}
{0.100 & & & & & & 0.050}

\[
V_{\ce{H2}}=0.050\times24.8=\SI{1.24}{\liter}.
\]

\textbf{Vậy,} thu được \SI{1.24}{\liter} \ce{H2}.
\end{loigiai}

\subsection{Điều chế ethylic alcohol}

Từ ethylene:
\[
\ce{C2H4 + H2O ->[H^+,t^\circ] C2H5OH}.
\]

Từ nguyên liệu chứa tinh bột, quá trình được mô tả khái quát:
\[
\text{tinh bột}\longrightarrow\text{glucose}\longrightarrow\text{ethylic alcohol}.
\]
Phản ứng lên men glucose:
\[
\ce{C6H12O6 ->[men] 2C2H5OH + 2CO2}.
\]

\begin{notechem}
Ethylic alcohol được dùng làm dung môi, nhiên liệu và nguyên liệu hóa học. Việc lạm dụng đồ uống có cồn gây nhiều tác hại đối với sức khỏe và an toàn xã hội.
\end{notechem}

\begin{tuluyen}
Một dung dịch rượu có thể tích \SI{750}{\milli\liter}, độ cồn $32^\circ$. Tính thể tích ethylic alcohol nguyên chất trong dung dịch.
\dapso{\SI{240}{\milli\liter}.}
\end{tuluyen}
